\section{Bison: Deklarasi, Aturan, dan Konflik}
\label{sec:ch13-bison}

\subsection{Struktur Berkas Bison}

Berkas Bison (\file{.y}) terdiri atas tiga bagian:
\begin{enumerate}[label=\arabic*.,leftmargin=*]
  \item \textbf{deklarasi}: token, tipe, precedence, dan kode C;
  \item \textbf{aturan}: grammar dan semantic action;
  \item \textbf{kode pengguna}: fungsi bantu dan \code{yyerror}.
\end{enumerate}

\subsection{Deklarasi}

Bagian deklarasi mendefinisikan token, tipe nilai, dan precedence.
Potongan \file{parser.y} ditunjukkan pada \lstref{lst:ch13-decl}.

\includecode[
  language=Bison,
  firstline=1,
  lastline=50,
  firstnumber=1,
  caption={Deklarasi \file{parser.y}},
  label={lst:ch13-decl}
]{../../code/subsetc/src/parser.y}

\code{\%union} mendefinisikan tipe nilai semantik. Setiap token
memiliki nilai bertipe \code{Token}, sedangkan nonterminal memiliki
nilai bertipe \code{AstNode *}. Deklarasi \code{\%type} menghubungkan
nonterminal dengan tipe tersebut.

\subsection{Aturan dan Aksi}

Aturan grammar ditulis dengan simbol \code{:} dan \code{|}. Setiap
alternatif dapat memiliki blok aksi C. Potongan aturan
\code{program} dan \code{stmt\_list} ditunjukkan pada
\lstref{lst:ch13-program}.

\includecode[
  language=Bison,
  firstline=53,
  lastline=82,
  firstnumber=53,
  caption={Aturan \code{program} dan \code{stmt\_list}},
  label={lst:ch13-program}
]{../../code/subsetc/src/parser.y}

Aksi pada \code{program} membuat simpul \code{AST\_PROGRAM} dan memindahkan
setiap pernyataan dari daftar sementara ke daftar anak simpul.
Pemindahan ini diperlukan karena Bison membangun daftar melalui
\code{next}, sedangkan AST menyimpan anak melalui \code{children}.

\subsection{Konflik}

Laporan konflik shift-reduce dan reduce-reduce dari Bison perlu
dibaca, bukan diabaikan \citep{thain2024}. Precedence dideklarasikan
dengan:
\begin{lstlisting}[language=Bison]
%left T_OR
%left T_AND
%left T_EQ T_NE
%left T_LT T_LE T_GT T_GE
%left T_PLUS T_MINUS
%left T_STAR T_SLASH T_PERCENT
%right T_NOT UMINUS
\end{lstlisting}
Deklarasi ini menyelesaikan konflik ekspresi tanpa mengubah grammar.
Untuk dangling else, Bison secara default memilih shift, yang sesuai
dengan aturan else terdekat.

\begin{warningbox}
Konflik yang tidak disengaja harus diperbaiki, bukan disembunyikan.
Periksa keluaran Bison dengan opsi \code{-v} untuk melihat
\file{parser.output} dan lokasi konflik.
\end{warningbox}
